\documentclass[12pt]{beamer}
\usepackage[utf8]{inputenc}
\usepackage[T1]{fontenc}
\usetheme{Copenhagen}
\usepackage[brazil]{babel}
\usepackage{fmtcount}
\usepackage{graphicx}
\usepackage{tikz}
\usepackage{booktabs}
\graphicspath{{figuras/}}
\usetikzlibrary{graphs}
\usepackage{ragged2e}



% setas para os diagramas
\tikzset{%
	>=latex, %melhor aparência para as setas
	inner sep=0pt,outer sep=0pt, %sem separação entre os nós
}


% numeracao dos slides
\expandafter
\def
\expandafter
\insertshorttitle
\expandafter{%
	\insertshorttitle
	\hfill%
		\insertframenumber
	\,/\,
	\inserttotalframenumber
}

% 
\begin{document}
	\author{Dalton Felipe, Yasmin Braga}
	\title{\large Introdução à Ciência dos dados com Python}
	\logo{
		\includegraphics[width=0.97 cm]{sc_logo}
		\includegraphics[width=1.17 cm]{labic}
		\includegraphics[width=0.8 cm]{brasao_ufopa}
	}
	%\subtitle{\large (Aula 01)}
	\institute{\large III Semana Cetec}
	%\date{}
	
	% slide principal
	\begin{frame}[plain]
		\centering
		\includegraphics[width=1.5cm]{figuras/brasao_ufopa.jpg}
		\\Universidade Federal do Oeste do Pará
		\\Instituto de Engenharia e Geociências
		\\Programa de Ciência e Tecnologia
		
		\maketitle
	\end{frame}

	% sumario
	\section[]{}
	\begin{frame}{Sumário}
		\tableofcontents
	\end{frame}
	
	% importa os slides do arquivo slides.tex
	
\section{Introdução}
	% slide 1
	\subsection{O que é ciência dos dados?}
		\begin{frame}{Introdução}
			
			\begin{block}{O que é Ciência dos dados?}
				Ciência dos dados pode ser definida como um processo de extrair insights de dados confusos \cite{grus:2015} . Ou seja, transformar dados ilegíveis em em informações poderosas acerca do passado e futuro. 
			\end{block}
			
		\end{frame}

\setbeamercovered{transparent}
\section{Justificativa}
\subsection{Porquê aprender Python?}
	\begin{frame}{Porquê aprender Python?}
		\begin{columns}
			
			\column{0.33\textwidth}
				\centering
				\begin{block}{\centering Fácil de aprender}
					\justifying
					\tiny Uma das maiores vantagens do Python é a forma intuitiva e fácil com que você consegue aprender a utilizá-la! Como sua sintaxe é extremamente simples, basta algumas noções de programação para começar a se aprofundar nessa linguagem!
				\end{block}
				
			\column{0.33\textwidth}
				\pause
				
				\begin{block}{\centering Multiplataforma}
					
					\justifying
					\tiny Por ser uma linguagem interpretada -- e não compilada para uma linguagem de máquina --, Python roda em diferentes plataformas. Além disso, a linguagem é multiparadigma, isto é, passível de desenvolvimento para qualquer ambiente e interface, como web, móvel ou desktop.
				\end{block}
			
			\column{0.33\textwidth}
				\pause
				\begin{block}{\centering Mercado de trabalho}
					
					\justifying
					\tiny Uma pesquisa mostra Python em primeiro lugar entre as linguagens mais adotadas — à frente das consolidadas Java e C++. Também demonstra que simplicidade não é sinônimo de limitação: Python é um instrumento poderoso de desenvolvimento de aplicações, aliando intuitividade e eficiência.
				\end{block}
					
		\end{columns}
	\end{frame}


\section{Objetivos}
\subsection{Quais os objetivos do curso?}
	\begin{frame}{Objetivos}
		
		\begin{block}{Quais os objetivos do curso?}
			\begin{itemize}[<+->]
				
				\item Oferecer um curso teórico e prático sobre a linguagem Python focada na análise de dados.
				
				\item Apresentar a sintaxe básica da linguagem para quem quer iniciar ou aprimorar o conhecimento na área de Data Science.
				
				\item Abordar algumas ferramentas e técnicas que auxiliam e simplificam a manipulação de dados. 
			\end{itemize}
			
		\end{block}
		
	\end{frame}

\section{Metodologia}
\subsection{Metodologia do curso}
	\begin{frame}{Metodologia}
		\begin{block}{Metodologia do curso}
			\begin{itemize}[<+->]
				\item O curso se dará na apresentação da linguagem Pyhton e sua sintaxe básica, bem como tipos de dados, estruturas de decisão, repetição dentre outras funções internas.
				
				\item Além disso serão aplicados exercícios práticos envolvendo o que foi apresentado.
			\end{itemize}
		\end{block}
	\end{frame}

\section{Conteúdo}
\begin{frame}{Conteúdo}

	\begin{columns}
		\column{0.5\linewidth}
			
			\begin{block}{Tipos de dados}
				
				\begin{itemize}
					\item strings 
					\item inteiros
					\item floats
					\item listas e dicts
				\end{itemize}
			\end{block}
			
			\begin{block}{Estruturas de controle}
				\begin{itemize}
					\item if 
					\item else
					\item elif
				\end{itemize}
			\end{block}
			
				
			
		\column{0.5\linewidth}
			
			\begin{block}{Estruturas de repetição}
				\begin{itemize}
					\item for
					\item while
				\end{itemize}
			\end{block}
		
			\begin{block}{Definição de Funções}
				\begin{itemize}
					\item def
				\end{itemize}
			\end{block}
		
			\begin{block}{Bibliotecas auxiliares}
				\begin{itemize}
					\item Numpy
					\item Matplotlib
					\item Pandas
				\end{itemize}
			\end{block}

	\end{columns}
	

\end{frame}



\section{Conclusão}

	
	\section{Referências}
	\begin{frame}{Referências}
		\bibliography{referencias}
		\bibliographystyle{plain}
	\end{frame}

	
	\begin{frame}
		\centering
		\huge{FIM}
	\end{frame}

\end{document}